Đề tài đã nghiên cứu và xây dựng hệ thống chatbot hỗ trợ tra cứu TTHC
cho tỉnh Lai Châu dựa trên kỹ thuật Retrieval-Augmented Generation. Hệ
thống kết hợp dữ liệu TTHC đã chuẩn hóa, vector database, mô hình
embedding, reranker và LLM để trả lời câu hỏi của người dùng theo hướng
tự nhiên nhưng vẫn bám sát dữ liệu nguồn.

Về mặt sản phẩm, đề tài đã triển khai backend FastAPI, frontend React,
Firebase Authentication, Firestore lưu lịch sử chat, ChromaDB lưu
vector, cơ chế intent nhanh, chuẩn hóa synonym, RAG pipeline, fallback
LLM và dashboard quản trị. Các chức năng chính như hỏi đáp chatbot, quản
lý phiên hội thoại, xác thực, quản lý tài liệu, đồng bộ vector database
và thống kê hệ thống đã được hiện thực hóa.

Về mặt ý nghĩa thực tiễn, hệ thống giúp giảm khó khăn của người dân khi
tra cứu thông tin hành chính, hỗ trợ đặt câu hỏi bằng ngôn ngữ tự nhiên,
giảm thời gian đọc văn bản dài và hỗ trợ cán bộ quản trị cập nhật dữ
liệu thủ tục linh hoạt hơn.

Trong tương lai, hệ thống có thể được phát triển theo các hướng: tích
hợp trực tiếp với cổng dịch vụ công để tra cứu trạng thái hồ sơ; bổ sung
chức năng hỏi đáp bằng giọng nói; xây dựng bộ đánh giá định lượng lớn
cho từng nhóm thủ tục; tối ưu tốc độ inference bằng cloud AI; bổ sung
chức năng trích dẫn nguồn chi tiết trong câu trả lời; và mở rộng sang
nhiều tỉnh, thành phố khác.
